\section{Exact distance coefficients for weighted norm trees}
\label{sec:tree-distance}

The finite-accuracy estimate above is uniform over objectives.  For a fixed
objective, one can determine the leading distance coefficient more precisely.
An internal subtree on which the objective vanishes contributes half its
barrier weight to the squared coefficient.  This contribution is invisible
to the rank of a nonzero exposing certificate, because the corresponding
dual block is zero.  It nevertheless cannot be avoided: the subtree must
shrink toward the vertex of its Lorentz cone.

We use the norm-tree domain $\Omega_T$ of
Section~\ref{sec:concrete-barriers}.  Write $V$ for its internal nodes,
$b=|V|$, and $w$ for its vector of leaf coordinates.  Fix weights
$\omega_v\geq1$ and the displayed barrier
\[
 F_\omega=-\sum_{v\in V}\omega_v\log q_v,\qquad
 W=\sum_{v\in V}\omega_v.
\]
It is a self-concordant barrier with positive definite restricted Hessian.
All distances in this section refer to this Hessian, with the root fixed
to one and all linking equations imposed.  Fix a unit vector $c$ in the
leaf space, put $g(w)=1-c^Tw$, and define
\[
 \alpha_v=\norm{c|_{\text{leaves below }v}},\qquad
 V_+=\{v:\alpha_v>0\},\quad V_0=V\setminus V_+,
 \qquad
 K_c=\sum_{v\in V_+}\omega_v+\frac12\sum_{v\in V_0}\omega_v.
\]
The root belongs to $V_+$; the set $V_0$ is a union of complete subtrees.

\begin{theorem}[Weighted norm-tree distance]\label{thm:tree-exact-distance}
For every fixed $x_0\in\Omega_T$, as $\epsilon\downarrow0$,
\begin{equation}\label{eq:tree-exact-distance}
 d_{F_\omega}(x_0,\{x\in\Omega_T:g(w)\leq\epsilon\})
 =\sqrt{K_c}\log(1/\epsilon)+O(\log\log(1/\epsilon)).
\end{equation}
More precisely, the lower bound has an $O(1)$ remainder; the displayed
$O(\log\log(1/\epsilon))$ is needed only for our upper construction.
The constants may depend on the fixed tree, weights, objective and reference.
\end{theorem}
\begin{proof}
We first construct a potential that bounds every feasible path.
For the primal Lorentz block $x_v=(t_v,y_v)$, define the Euclidean dual vector
\[
 s_v=(\alpha_v,-(\alpha_u)_{u\in I(v)},-(c_j)_{j\in J(v)}).
\]
Here $I(v)$ and $J(v)$ are the internal-child and leaf-child sets.
The identity $\alpha_v^2=\sum_{u\in I(v)}\alpha_u^2+
\sum_{j\in J(v)}c_j^2$ says that $s_v$ is a nonzero boundary dual
vector for $v\in V_+$, and is zero otherwise.  Consequently
\[
 \beta_v=\ip{s_v}{x_v}>0\quad(v\in V_+),\qquad
 \sum_{v\in V_+}\beta_v=1-c^Tw=g(w).
\]
The sum telescopes, including all internal links.  Thus each
$\beta_v\leq g$.

There is also a uniform bound $q_z\leq C_c g$ for every inactive node
$z\in V_0$.  To see this, first let $u\in V_0$ have an active parent
$v$.  Write $e_v=((\alpha_u)_{u\in I(v)},(c_j)_{j\in J(v)})/\alpha_v$,
which is a unit vector.  The coordinate occupied by $t_u$ is orthogonal
to $e_v$.  Since $\norm{y_v}<t_v$ and $t_v\leq1$,
\[
 t_u^2\leq\norm{y_v-(e_v^Ty_v)e_v}^2
 \leq t_v^2-(e_v^Ty_v)^2
 \leq 2t_v(t_v-e_v^Ty_v)
 =\frac{2t_v\beta_v}{\alpha_v}\leq\frac{2g}{\alpha_v}.
\]
Every descendant axis is at most $t_u$, so its determinant is at most
$t_u^2$.  Taking the maximum of $2/\alpha_v$ over the finitely many
active parents of inactive subtrees proves the assertion.

For clarity, the needed dual-norm identities can be checked without
Jordan normalization.  For $f(x)=-\log(x^TJx)$ on a Lorentz cone,
$J=\operatorname{Diag}(1,-I)$ and $q=x^TJx$, direct inversion gives
\[
 (D^2f(x))^{-1}=xx^T-\frac q2J.
\]
If $s\ne0$ is a boundary dual vector, then $s^TJs=0$ and
$\norm{D\log(s^Tx)}_{(D^2f)^{-1}}^2=1$.
Logarithmic homogeneity likewise gives
$\norm{Df}_{(D^2f)^{-1}}^2=2$.
It follows that the potential
\[
 \Psi(x)=-\sum_{v\in V_+}\omega_v\log\beta_v
          -\frac12\sum_{v\in V_0}\omega_v\log q_v
\]
has squared dual gradient norm exactly $K_c$ in the ambient product
metric $D^2F_\omega$.  Restriction to the linked root slice can only
decrease this norm: its value is the supremum over fewer tangent directions.
At every point with $g\leq\epsilon$, the preceding determinant and
pairing bounds imply
\[
 \Psi(x)-\Psi(x_0)\geq K_c\log(1/\epsilon)-O(1).
\]
Integrating $|D\Psi[h]|\leq\sqrt{K_c}\norm{h}_{F_\omega,x}$
along an arbitrary path proves the claimed lower bound.

For the upper bound put $A=|V_+|$, $Z=|V_0|$, and, for large
$\lambda$, set $h=e^{-\lambda}$ and prescribe
\begin{equation}\label{eq:tree-distance-curve}
 q_v=\begin{cases}h,&v\in V_+,\\h/\lambda^2,&v\in V_0,\end{cases}
 \quad r^2=1-Ah-Zh/\lambda^2,
 \quad w=rc,\quad
 t_v^2=r^2\alpha_v^2+\sum_{z\in V\text{ below }v}q_z.
\end{equation}
The last sum includes $v$.  Choose the positive square roots.  The subtree
identities give precisely the prescribed determinants, and the root
equation gives $t_{\mathrm{root}}=1$.  Thus this is an interior feasible curve.
Its gap is $1-r=Ah/2+O(h/\lambda^2+h^2)$, since $A\geq1$.

A dot will denote differentiation with respect to $\lambda$.
The Lorentz metric has the identity
\begin{equation}\label{eq:tree-speed-identity}
 D^2f(x)[\dot x,\dot x]
 =\left(\frac{\dot q}{q}\right)^2
       -\frac{2(\dot t^2-\norm{\dot y}^2)}q.
\end{equation}
At an active node, $t_v$ tends to $\alpha_v>0$ and
$\dot t_v=O(h)$.  The same holds for active child axes and for
leaf-coordinate derivatives; inactive child derivatives are
$O(\sqrt h/\lambda)$.  Since $q_v=h$, this block contributes
$1+O(h+\lambda^{-2})$ to the unweighted squared speed.
At an inactive node, all its leaf coordinates vanish and all its
descendants are inactive.  Its whole block is a fixed interior Lorentz
vector multiplied by $\sqrt h/\lambda$.  Its squared speed is therefore
exactly $2(1/2+1/\lambda)^2$.  Adding the weighted contributions gives
\[
 \norm{\dot x}_{F_\omega,x}^2=K_c+O(1/\lambda).
\]
Starting the curve at any sufficiently large fixed $\lambda_0$, its
length to $\lambda$ is $\sqrt{K_c}\lambda+O(\log\lambda)$.
Connect $x_0$ to the starting point by an interior line segment, whose
length is finite.  Along the curve,
$\lambda=\log(1/g)+O(1)$, which proves the upper bound.
\end{proof}

\begin{corollary}[Central-path excess and independent products]
\label{cor:tree-central-excess}
For the weighted tree, the central route from its analytic center to
gap $\epsilon$ has length
$\sqrt W\log(1/\epsilon)+O(1)$.  Its ratio to the shortest distance
to that target tends to $\sqrt{W/K_c}$.  It is asymptotically optimal
in leading coefficient exactly when every internal subtree has nonzero
objective support.

For independent trees and the objective
$\sum_a\lambda_a c_a^Tw_a$, with fixed $\lambda_a>0$ and unit $c_a$,
the squared distance coefficient is $\sum_a K_{c_a}$.
Independent factors with zero objective weight contribute zero.
\end{corollary}
\begin{proof}
The weighted central-path formulas of the preceding section give the first
length.  The ratio follows from Theorem~\ref{thm:tree-exact-distance};
$K_c=W$ holds exactly when $V_0$ is empty.
For products, sum the lower potentials.  Each individual gap is at most
$\epsilon/\lambda_a$, which changes only the fixed remainder, and the
squared dual norms add.  For the upper bound synchronize the parameter
$\lambda$ in \eqref{eq:tree-distance-curve} across all positively
weighted trees.  The total objective gap is a fixed positive constant
times $e^{-\lambda}(1+o(1))$, and squared speeds add.  Keep the remaining
independent factors at their reference points.
\end{proof}

For example, take a tree with two internal nodes, with the objective
supported entirely on a root leaf.  For unit weights, the root is active
and its internal child is inactive.  The shortest-distance coefficient
is $\sqrt{3/2}$, whereas the central coefficient is $\sqrt2$.
If the objective also has nonzero support in the child subtree, both
coefficients are $\sqrt2$.  Thus the coefficient need not vary continuously
as a component of the objective tends to zero: the asymptotic here holds
with the objective fixed before accuracy tends to zero.

The classical study of barrier metrics and central-path distances
\citep{NT2002,NesterovNemirovski2008}, and the uniform norm-tree estimates
in \citet{CentralPathCompanion}, provide the context for this calculation.
The specific result here is the explicit weighted coefficient $K_c$,
including objective-free subtrees, for this fixed family of linked
Lorentz barriers.  It does not identify the best barrier on the tree
domain, nor assert an iteration lower bound without a restriction on
the metric distance traversed per accepted step.
